\section{Jeux palisade}

Le problème \textit{Palisade} présente plusieurs points communs avec le 
problème précédent de \textit{Loopy}, mais sa structure est plus riche. 
On considère à nouveau un ensemble de cases formant une grille rectangulaire. 
Cette fois, l'objectif n'est plus de construire une boucle, mais de partitionner la grille en régions connexes, 
disjointes deux à deux et sans superposition.

Chaque région possède une taille fixée, notée \(s\), correspondant au nombre de cases qui la composent. 
Dans le cas général du problème \textit{Palisade}, on prend \(s = 5\). En parallèle, certaines cases portent un 
indice numérique, comme dans \textit{Loopy}. Cet indice indique combien de côtés de la case considérée appartiennent 
à une frontière entre deux régions. Autrement dit, pour chaque case numérotée, le nombre inscrit correspond au nombre 
d'arêtes de son contour qui délimitent une région de \textit{Palisade}. Comme une case carrée possède quatre côtés, 
cet indice est nécessairement compris entre \(0\) et \(4\), pour un $s > 1$.

\subsection{Modélisation}

\subsubsection{Notations}

On introduit les ensembles et paramètres suivants :

\begin{itemize}
    \item \(\mathcal{C}\) : ensemble des cases de la grille ;
    \item \(\mathcal{A}\) : ensemble des paires de cases adjacentes par un côté ;
    \item \(s\) : taille imposée de chaque région, avec \(s = 5\) dans le cas général ;
    \item \(K = \card{\mathcal{C}}/s\) : nombre total de régions ;
    \item \(\mathcal{P} = \{1, \dots, K\}\) : ensemble des identifiants de régions ;
    \item \(\mathcal{Q} \subseteq \mathcal{C}\) : ensemble des cases contenant un indice ;
    \item pour toute case \(c \in \mathcal{Q}\), \(a_c\) désigne la valeur inscrite dans la case \(c\) ;
    \item pour toute case \(c \in \mathcal{C}\), \(\delta_c\) désigne le nombre de côtés de la case \(c\) situés sur le bord extérieur de la grille.
\end{itemize}

Les bords extérieurs de la grille sont toujours considérés comme des palissades.

\subsubsection{Variables de décision}

On introduit plusieurs familles de variables binaires.

\paragraph{Affectation des cases aux régions}
\mbox{}\\

Pour chaque case \(c \in \mathcal{C}\) et pour chaque région \(p \in \mathcal{P}\), on définit :
\[
x_{c,p} \in \binary
\]
telle que :
\[
x_{c,p} = 1 \iff \text{la case } c \text{ appartient à la région } p.
\]

\paragraph{Présence d'une palissade entre deux cases voisines}
\mbox{}\\

Pour chaque paire \(\{u,v\} \in \mathcal{A}\), on définit :
\[
y_{uv} \in \binary
\]
telle que :
\[
y_{uv} = 1 \iff \text{une palissade sépare les cases } u \text{ et } v.
\]

Il est également nécessaire d'introduire une variable auxiliaire. Même si elle ne fait pas directement partie de la solution recherchée, elle permet d'écrire linéairement certaines relations entre les variables précédentes.

\paragraph{Variables auxiliaires pour détecter si deux cases voisines sont dans la même région}
\mbox{}\\

Pour chaque paire \(\{u,v\} \in \mathcal{A}\) et pour chaque région \(p \in \mathcal{P}\), on introduit :
\[
z_{uv,p} \in \binary
\]
telle que :
\[
z_{uv,p} = 1 \iff u \text{ et } v \text{ sont adjacentes et appartiennent toutes deux à la région } p.
\]


\subsubsection{Contraintes}

Les deux premières contraintes sont des contraintes de base. La première traduit le fait que chaque case doit appartenir à une unique région. La seconde compte le nombre de cases affectées à chaque région ; cette somme doit alors être égale à la taille imposée \(s\).

\paragraph{Chaque case appartient à une seule région}
\mbox{}\\

Pour toute case \(c \in \mathcal{C}\), on impose :
\[
\sum_{p \in \mathcal{P}} x_{c,p} = 1
\qquad \forall c \in \mathcal{C}
\]

\paragraph{Chaque région contient exactement \(s\) cases}
\mbox{}\\

Pour toute région \(p \in \mathcal{P}\), on impose :
\[
\sum_{c \in \mathcal{C}} x_{c,p} = s
\qquad \forall p \in \mathcal{P}
\]

\paragraph{Linéarisation de l'appartenance simultanée à une même région}
\mbox{}\\

Pour chaque paire adjacente \(\{u,v\} \in \mathcal{A}\) et pour chaque région \(p \in \mathcal{P}\), on impose :

\[
z_{uv,p} \leq x_{u,p}
\qquad \forall \{u,v\} \in \mathcal{A},\; \forall p \in \mathcal{P}
\]

Cette première contrainte signifie que si la case \(u\) n'appartient pas à la région \(p\), alors la variable \(z_{uv,p}\) ne peut pas valoir \(1\). Autrement dit, la paire de cases adjacentes \(\{u,v\}\) ne peut pas être comptée dans la région \(p\) si \(u\) n'y appartient pas.

\[
z_{uv,p} \leq x_{v,p}
\qquad \forall \{u,v\} \in \mathcal{A},\; \forall p \in \mathcal{P}
\]

De la même manière, cette deuxième contrainte exprime le même raisonnement du point de vue de la case \(v\). Si \(v\) n'appartient pas à la région \(p\), alors \(z_{uv,p}\) doit être nul.

\[
z_{uv,p} \geq x_{u,p} + x_{v,p} - 1
\qquad \forall \{u,v\} \in \mathcal{A},\; \forall p \in \mathcal{P}
\]

Enfin, cette troisième contrainte force \(z_{uv,p}\) à valoir \(1\) lorsque les deux cases \(u\) et \(v\) appartiennent simultanément à la région \(p\). En effet, si aucune des deux cases n'appartient à \(p\), alors le membre de droite vaut \(-1\), ce qui ne crée aucune contrainte supplémentaire. Si une seule des deux cases appartient à \(p\), alors le membre de droite vaut \(0\), et \(z_{uv,p}\) peut seulement rester nul. En revanche, si les deux cases appartiennent à la région \(p\), alors le membre de droite vaut \(1\), ce qui impose \(z_{uv,p} = 1\).

Ainsi, \(z_{uv,p} = 1\) si et seulement si les deux cases adjacentes \(u\) et \(v\) appartiennent à la même région \(p\).

\paragraph{Définition des palissades}
\mbox{}\\

Cette contrainte est directement liée à la modélisation précédente. Deux cas peuvent alors se présenter. Si les deux cases adjacentes appartiennent à la même région, il ne doit pas y avoir de palissade entre elles. En revanche, si les deux cases appartiennent à des régions différentes, alors une palissade doit nécessairement les séparer.

On modélise cela par :
\[
y_{uv} + \sum_{p \in \mathcal{P}} z_{uv,p} = 1
\qquad \forall \{u,v\} \in \mathcal{A}
\]

En effet :
\begin{itemize}
    \item si \(u\) et \(v\) sont dans la même région, alors \(\sum_{p \in \mathcal{P}} z_{uv,p} = 1\), donc \(y_{uv} = 0\) ;
    \item sinon, \(\sum_{p \in \mathcal{P}} z_{uv,p} = 0\), donc \(y_{uv} = 1\).
\end{itemize}

\paragraph{Contraintes liées aux indices}
\mbox{}\\

Pour chaque case numérotée \(c \in \mathcal{Q}\), le nombre total de côtés portant une palissade doit être égal à l'indice inscrit.

Si \(N(c)\) désigne l'ensemble des cases voisines de \(c\), alors :
\[
\delta_c + \sum_{v \in N(c)} y_{cv} = a_c
\qquad \forall c \in \mathcal{Q}
\]

où :
\begin{itemize}
    \item \(\delta_c\) compte les côtés extérieurs de la grille, qui sont toujours des palissades ;
    \item \(\sum_{v \in N(c)} y_{cv}\) compte les palissades internes autour de la case \(c\).
\end{itemize}

\paragraph{Contrainte difficile : connexité des régions}
\mbox{}\\

La difficulté principale du problème consiste à imposer que chaque région soit connexe. Pour modéliser cela, on suppose qu'un flot doit circuler dans chaque région à partir d'une case racine.

\par
\begin{adjustwidth}{1.5em}{}

\noindent\textbf{Choix d'une racine par région}
\par

Pour chaque case \(c \in \mathcal{C}\) et chaque région \(p \in \mathcal{P}\), on introduit :
\[
r_{c,p} \in \binary
\]
avec :
\[
r_{c,p} = 1 \iff \text{la case } c \text{ est choisie comme racine de la région } p.
\]

On impose ensuite :
\[
\sum_{c \in \mathcal{C}} r_{c,p} = 1
\qquad \forall p \in \mathcal{P}
\]

\[
r_{c,p} \leq x_{c,p}
\qquad \forall c \in \mathcal{C},\; \forall p \in \mathcal{P}
\]

Ainsi, chaque région possède exactement une racine, qui doit appartenir à cette région. La seconde contrainte signifie donc qu'une case ne peut être choisie comme racine de la région \(p\) que si elle appartient effectivement à cette région.

\par
\noindent\textbf{Variables de flot}
\par

Pour chaque région \(p \in \mathcal{P}\) et pour chaque arc orienté \((u,v)\) associé à une adjacency \(\{u,v\} \in \mathcal{A}\), on introduit une variable continue :
\[
f^{p}_{uv} \geq 0
\]

Cette variable représente la quantité de flot envoyée de \(u\) vers \(v\) dans la région \(p\).

\par
\noindent\textbf{Le flot ne circule qu'à l'intérieur d'une région}
\par

On impose :
\[
f^{p}_{uv} \leq (s-1)x_{u,p}
\qquad \forall (u,v),\; \forall p \in \mathcal{P}
\]
\[
f^{p}_{uv} \leq (s-1)x_{v,p}
\qquad \forall (u,v),\; \forall p \in \mathcal{P}
\]

Dans le cas général d'exemple considéré ici, on prend \(s = 5\). Ces contraintes deviennent donc :
\[
f^{p}_{uv} \leq 4x_{u,p}
\qquad \forall (u,v),\; \forall p \in \mathcal{P}
\]
\[
f^{p}_{uv} \leq 4x_{v,p}
\qquad \forall (u,v),\; \forall p \in \mathcal{P}
\]

Ces inégalités signifient que si l'arc \((u,v)\) appartient à la région \(p\), alors le flot qui y circule peut prendre une valeur comprise entre \(1\) et \(s-1\), selon la quantité de flot qu'il reste à transmettre. À l'intérieur d'une région de taille \(s\), la racine peut en effet envoyer au plus \(s-1\) unités, puis cette quantité décroît progressivement jusqu'à la dernière unité de flot.

On rappelle également que sur un arc correspondant à une palissade, le flot est nécessairement nul. En effet, si \(u\) et \(v\) appartiennent à deux régions différentes, alors pour une région \(p\) fixée, l'une des deux variables \(x_{u,p}\) ou \(x_{v,p}\) vaut \(0\). L'une des deux contraintes précédentes impose donc \(f^{p}_{uv} \leq 0\). Comme on a déjà défini \(f^{p}_{uv} \geq 0\), on obtient finalement \(f^{p}_{uv} = 0\).

\par
\noindent\textbf{Conservation du flot}
\par

Pour finir, il faut imposer que le flot se conserve à l'intérieur de chaque région. Pour chaque case \(c \in \mathcal{C}\) et chaque région \(p \in \mathcal{P}\), on impose :
\[
\sum_{u:(u,c)} f^{p}_{uc} - \sum_{v:(c,v)} f^{p}_{cv} = x_{c,p} - s\,r_{c,p}
\qquad \forall c \in \mathcal{C},\; \forall p \in \mathcal{P}
\]

Dans le cas général d'exemple considéré ici, avec \(s = 5\), on obtient :
\[
\sum_{u:(u,c)} f^{p}_{uc} - \sum_{v:(c,v)} f^{p}_{cv} = x_{c,p} - 5r_{c,p}
\qquad \forall c \in \mathcal{C},\; \forall p \in \mathcal{P}
\]
L'interprétation que l'on peut en faire est la suivante :

\begin{itemize}
    \item Si la case \(c\) appartient à la région \(p\) et n'est pas la racine, alors \(x_{c,p} = 1\) et \(r_{c,p} = 0\). Le membre de droite vaut donc \(1\), ce qui donne :
    \[
    \text{flux entrant} - \text{flux sortant} = 1
    \]
    La case doit donc recevoir une unité de flot.
    \item Si la case \(c\) n'appartient pas à la région \(p\), alors \(x_{c,p} = 0\) et \(r_{c,p} = 0\). Le membre de droite vaut donc \(0\). Il n'y a alors aucun flot à conserver, et le bilan entre le flux entrant et le flux sortant est nul.
    \item Si la case \(c\) appartient à la région \(p\) et qu'elle est la racine, alors \(x_{c,p} = 1\) et \(r_{c,p} = 1\). Le membre de droite vaut donc \(1-s\). Avec \(s=5\), on obtient :
    \[
    \text{flux entrant} - \text{flux sortant} = 1 - 5 = -4
    \]
    La racine doit donc envoyer \(s-1\) unités de flot aux autres cases de la région. Dans le cas \(s=5\), elle doit envoyer \(4\) unités.
\end{itemize}

Cette formulation garantit que toutes les cases d'une même région sont atteignables depuis la racine, donc que la région est connexe.

\end{adjustwidth}

\subsubsection{Fonction objectif}

Comme dans le problème \textit{Loopy}, il s'agit uniquement de trouver une solution réalisable. On choisit donc une fonction objectif constante :
\[
\min 0
\]

\subsubsection{Modèle final}

Le modèle complet s'écrit donc :
\[
\min 0
\]

sous les contraintes :

\begin{center}
\renewcommand{\arraystretch}{1.4}
\begin{tabularx}{\textwidth}{>{$}X<{$} l}
\sum_{p \in \mathcal{P}} x_{c,p} = 1 \qquad \forall c \in \mathcal{C} & 1 case dans 1 région \\
\sum_{c \in \mathcal{C}} x_{c,p} = s \qquad \forall p \in \mathcal{P} & taille des régions \\
z_{uv,p} \leq x_{u,p} \qquad \forall \{u,v\} \in \mathcal{A},\; \forall p \in \mathcal{P} & linéarisation sur \(u\) \\
z_{uv,p} \leq x_{v,p} \qquad \forall \{u,v\} \in \mathcal{A},\; \forall p \in \mathcal{P} & linéarisation sur \(v\) \\
z_{uv,p} \geq x_{u,p} + x_{v,p} - 1 \qquad \forall \{u,v\} \in \mathcal{A},\; \forall p \in \mathcal{P} & appartenance des voisins \\
y_{uv} + \sum_{p \in \mathcal{P}} z_{uv,p} = 1 \qquad \forall \{u,v\} \in \mathcal{A} & définition des palissades \\
\delta_c + \sum_{v \in N(c)} y_{cv} = a_c \qquad \forall c \in \mathcal{Q} & respect des indices de la case \\
\sum_{c \in \mathcal{C}} r_{c,p} = 1 \qquad \forall p \in \mathcal{P} & choix d'une racine \\
r_{c,p} \leq x_{c,p} \qquad \forall c \in \mathcal{C},\; \forall p \in \mathcal{P} & racine valide \\
f^{p}_{uv} \leq (s-1)x_{u,p} \qquad \forall (u,v),\; \forall p \in \mathcal{P} & capacité sur \(u\) pour l'arête\\
f^{p}_{uv} \leq (s-1)x_{v,p} \qquad \forall (u,v),\; \forall p \in \mathcal{P} & capacité sur \(v\) pour l'arête\\
\sum_{u:(u,c)} f^{p}_{uc} - \sum_{v:(c,v)} f^{p}_{cv} = x_{c,p} - s\,r_{c,p} \qquad \forall c \in \mathcal{C},\; \forall p \in \mathcal{P} & conservation du flot
\end{tabularx}
\end{center}

avec :
\[
x_{c,p},\; y_{uv},\; z_{uv,p},\; r_{c,p} \in \binary,
\qquad
f^{p}_{uv} \geq 0
\]

Cette formulation permet donc de modéliser le jeu \textit{Palisade} comme un programme linéaire en nombres entiers.

\subsection{Traitement de la règle difficile par callback}

\subsubsection{Pourquoi utiliser un callback ?}

La règle difficile du problème \textit{Palisade} est la connexité des régions. La formulation par flot présentée plus haut permet de l'imposer exactement, mais elle alourdit fortement le modèle lorsque la grille devient grande.

Posons \(N = \card{\mathcal{C}}\) le nombre total de cases et \(K = N/s\) le nombre de régions. La partie dédiée à la connexité ajoute alors :
\begin{itemize}
    \item \(\card{\mathcal{C}} \times K = NK\) variables binaires \(r_{c,p}\) pour le choix des racines ;
    \item \(2\card{\mathcal{A}} \times K\) variables continues \(f^p_{uv}\) pour les flots sur les arcs orientés ;
    \item \(K + NK + 4\card{\mathcal{A}}K + NK\) contraintes supplémentaires pour le choix des racines, les bornes sur le flot et la conservation du flot.
\end{itemize}

Ici, \(\mathcal{A}\) désigne l'ensemble des paires de cases adjacentes par un côté. Ainsi, \(\card{\mathcal{A}}\) compte le nombre total de voisinages dans la grille, c'est-à-dire le nombre de paires \(\{u,v\}\) de cases qui partagent une arête.

Le terme
\[
K + NK + 4\card{\mathcal{A}}K + NK
\]
provient alors directement du détail des contraintes de connexité :
\begin{itemize}
    \item \(K\) contraintes pour imposer qu'il y ait exactement une racine par région :
    \[
    \sum_{c \in \mathcal{C}} r_{c,p} = 1
    \qquad \forall p \in \mathcal{P} ;
    \]
    \item \(NK\) contraintes pour imposer qu'une racine appartienne bien à sa région :
    \[
    r_{c,p} \leq x_{c,p}
    \qquad \forall c \in \mathcal{C},\; \forall p \in \mathcal{P} ;
    \]
    \item \(4\card{\mathcal{A}}K\) contraintes pour borner le flot sur les arcs orientés : pour chaque paire adjacente, on considère deux orientations possibles \((u,v)\) et \((v,u)\), et pour chacune de ces orientations on écrit deux bornes, l'une avec \(x_{u,p}\) et l'autre avec \(x_{v,p}\) ;
    \item enfin, \(NK\) contraintes de conservation du flot, une pour chaque case \(c\) et chaque région \(p\).
\end{itemize}

Or, dans une grille rectangulaire, on a \(\card{\mathcal{A}} = O(N)\). En effet, si la grille possède \(n\) lignes et \(m\) colonnes, alors \(N = nm\), tandis que le nombre de paires adjacentes vaut
\[
n(m-1) + m(n-1) = 2nm - n - m.
\]
Cette quantité est donc linéaire en \(nm\), c'est-à-dire en \(N\). Comme \(K = N/s\) et que \(s\) est fixé dans le problème, le nombre de variables et de contraintes supplémentaires est de l'ordre de
\[
O\big((\card{\mathcal{A}} + \card{\mathcal{C}})K\big) = O(N^2).
\]

Autrement dit, la taille du modèle croît quadratiquement avec le nombre de cases dès que l'on ajoute la connexité par flot. Cette augmentation se répercute à la fois sur l'espace mémoire nécessaire et sur le temps de résolution, car chaque n\oe ud du \textit{branch-and-bound} doit manipuler une relaxation linéaire plus grande.

Pour de petites ou moyennes instances, cette formulation reste tout à fait exploitable. En revanche, pour des grilles plus grandes, il est souvent plus efficace de conserver un modèle de base plus léger et de traiter la connexité dynamiquement par \textit{callback}.

\subsubsection{Principe du callback}

L'idée est de ne pas introduire dès le départ toutes les contraintes de flot. On résout d'abord le modèle de base contenant les contraintes d'affectation, de taille des régions, de palissades et d'indices. Lorsqu'une solution entière candidate est trouvée, on vérifie alors si chaque région est connexe.

Pour une région fixée \(p\), on considère l'ensemble des cases telles que \(x_{c,p} = 1\), puis on reconstruit le sous-graphe d'adjacence correspondant. Si ce sous-graphe n'est pas connexe, on choisit une composante connexe stricte non vide \(W \subset \mathcal{C}\) de cette région et on ajoute une contrainte qui interdit cette composante isolée.

La coupe suivante permet de le faire :
\[
\sum_{c \in W} x_{c,p}
\leq
\card{W} - 1 + \sum_{\{u,v\} \in B(W)} z_{uv,p},
\]
où \(B(W)\) désigne l'ensemble des paires de cases adjacentes ayant une extrémité dans \(W\) et l'autre dans \(\mathcal{C} \setminus W\). Autrement dit, \(B(W)\) représente le bord de \(W\).

Le membre de gauche
\[
\sum_{c \in W} x_{c,p}
\]
compte combien de cases de \(W\) sont affectées à la région \(p\). Si toutes les cases de \(W\) appartiennent à \(p\), alors cette somme vaut exactement \(\card{W}\).

Le dernier terme du membre de droite
\[
\sum_{\{u,v\} \in B(W)} z_{uv,p}
\]
compte les connexions entre \(W\) et l'extérieur de \(W\) qui restent dans la même région \(p\). En effet, \(z_{uv,p}=1\) signifie que les deux cases voisines \(u\) et \(v\) appartiennent toutes les deux à la région \(p\).

Si \(W\) est isolé dans la région \(p\), aucune paire du bord \(B(W)\) ne relie \(W\) à une autre case de la même région. On a donc :
\[
\sum_{\{u,v\} \in B(W)} z_{uv,p} = 0,
\]
et la coupe devient :
\[
\sum_{c \in W} x_{c,p} \leq \card{W} - 1,
\]
ce qui interdit que toutes les cases de \(W\) restent affectées à la région \(p\). Le terme \(-1\) est donc essentiel : si on écrivait seulement \(\card{W}\), la contrainte serait toujours vraie et ne couperait pas la solution déconnectée.

En revanche, si au moins une paire du bord relie \(W\) à une autre case de la même région, alors
\[
\sum_{\{u,v\} \in B(W)} z_{uv,p} \geq 1.
\]
La coupe autorise alors \(\sum_{c \in W} x_{c,p} = \card{W}\), car \(W\) n'est plus une composante isolée.


Le \textit{callback} consiste donc à répéter les étapes suivantes :
\begin{enumerate}
    \item résoudre le modèle de base ;
    \item détecter les régions déconnectées dans la solution entière candidate ;
    \item ajouter une coupe de connexité pour chaque composante interdite ;
    \item relancer la recherche jusqu'à obtenir une solution où toutes les régions sont connexes.
\end{enumerate}

Il ne s'agit pas d'une heuristique, car on n'accepte jamais une solution déconnectée : chaque solution candidate qui viole la connexité est coupée par une contrainte valide. Le callback reste donc une méthode exacte, mais les contraintes de connexité ne sont ajoutées que lorsqu'elles deviennent nécessaires, ce qui permet de garder un modèle initial plus compact.
